\documentclass[11pt,a4paper]{article}
\usepackage[utf8]{inputenc}
\usepackage[spanish]{babel}
\usepackage{amsmath,amssymb,amsfonts}
\usepackage{algorithm}
\usepackage{algpseudocode}
\usepackage{geometry}
\usepackage{enumitem}
\usepackage{hyperref}
\usepackage{xcolor}
\usepackage{tcolorbox}
\usepackage{booktabs}

\geometry{margin=2.5cm}

\title{Guía de Respuestas Correctas\\Taller 1: Invariante de Bucle\\ECIN-608 -- Diseño y Análisis de Algoritmos}
\author{Guía de Estudio Académica}
\date{2do Semestre, 2026}

\newcommand{\inv}{I}

\begin{document}
\maketitle
\tableofcontents
\newpage

% ============================================================
% INSTRUCCIONES GENERALES
% ============================================================
\section{Instrucciones Generales para Obtener los 100 Puntos}

Para cada algoritmo, la presentación debe demostrar \textbf{competencia} en las siguientes áreas según la rúbrica:

\begin{itemize}
    \item \textbf{Contenido técnico (30 pts):} Algoritmo correcto, invariante formal, inicialización, mantenimiento, terminación, complejidad y casos límite.
    \item \textbf{Estructura y apoyo visual (20 pts):} 8--10 diapositivas, legibles, con diagramas/animaciones.
    \item \textbf{Comunicación oral (20 pts):} Explicación clara, vocabulario preciso, no leer diapositivas.
    \item \textbf{Preguntas y dominio (20 pts):} Respuestas fundamentadas, reconocer limitaciones.
    \item \textbf{Tiempo y participación (10 pts):} 5--7 minutos, ambos participantes activos.
\end{itemize}

\begin{tcolorbox}[colback=blue!5!white, colframe=blue!75!black, title={Estructura Obligatoria para Cada Algoritmo}]
\begin{enumerate}
    \item \textbf{Pseudocódigo} del algoritmo completo.
    \item \textbf{Definición formal del invariante} $I(i)$ con notación matemática.
    \item \textbf{Demostración de inicialización}: $I(i_0)$ verdadero.
    \item \textbf{Demostración de mantenimiento}: $I(i) \Rightarrow I(i+1)$.
    \item \textbf{Demostración de terminación}: $I(i_{\text{final}}) \wedge \neg\text{cond\_bucle} \Rightarrow \text{postcondición}$.
    \item \textbf{Análisis de complejidad}: temporal (mejor, peor, promedio) y espacial.
    \item \textbf{Casos límite}: discusión de al menos 4 casos especiales.
\end{enumerate}
\end{tcolorbox}

\newpage

% ============================================================
% ALGORITMO 1: BÚSQUEDA LINEAL
% ============================================================
\section{Algoritmo 1: Búsqueda Lineal}

\subsection{Pseudocódigo}
\begin{algorithm}[H]
\caption{Búsqueda Lineal}
\begin{algorithmic}[1]
\Require Arreglo $A[1..n]$, elemento objetivo $x$
\Ensure Índice $j$ tal que $A[j] = x$, o $-1$ si no existe
\State $j \gets -1$
\For{$i \gets 1$ \textbf{to} $n$}
    \If{$A[i] = x$}
        \State $j \gets i$
        \State \textbf{return} $j$
    \EndIf
\EndFor
\State \textbf{return} $j$
\end{algorithmic}
\end{algorithm}

\subsection{Invariante de Bucle}
Sea el bucle \texttt{for} que recorre $i$ desde $1$ hasta $n$. El invariante es:
\[
\boxed{\inv(i) \equiv \text{``No existe } k \in \{1, \ldots, i-1\} \text{ tal que } A[k] = x\text{''}}
\]
Es decir, antes de examinar $A[i]$, se ha verificado que $x$ no aparece en las posiciones $1$ a $i-1$.

\subsection{Demostración}

\subsubsection{Inicialización}
Antes de la primera iteración, $i = 1$. El conjunto $\{1, \ldots, 0\}$ es vacío, por lo que la afirmación ``no existe $k \in \emptyset$ tal que $A[k] = x$'' es \textbf{vacuamente verdadera}. $\square$

\subsubsection{Mantenimiento}
Supongamos $\inv(i)$ verdadero al inicio de la iteración $i$:
\begin{itemize}
    \item Si $A[i] = x$: el algoritmo retorna $i$ y el bucle termina. No hay iteración siguiente, por lo que la propiedad se mantiene trivialmente.
    \item Si $A[i] \neq x$: al inicio de la iteración $i+1$, el conjunto $\{1, \ldots, i\}$ no contiene $x$ (pues $\inv(i)$ afirma que no está en $\{1,\ldots,i-1\}$ y acabamos de verificar $A[i] \neq x$). Luego $\inv(i+1)$ es verdadero. $\square$
\end{itemize}

\subsubsection{Terminación}
El bucle termina en dos casos:
\begin{enumerate}
    \item \textbf{Retorno anticipado} ($A[i] = x$): Se retorna $i$ con $A[i] = x$, lo cual es correcto.
    \item \textbf{Bucle completa todas las iteraciones}: Al finalizar, $i = n+1$. El invariante $\inv(n+1)$ indica que $x \notin A[1..n]$. Se retorna $-1$, lo cual es correcto. $\square$
\end{enumerate}

\subsection{Complejidad}
\begin{itemize}
    \item \textbf{Mejor caso}: $O(1)$ --- el objetivo está en la primera posición.
    \item \textbf{Peor caso}: $O(n)$ --- el objetivo está en la última posición o no existe.
    \item \textbf{Promedio}: $O(n)$ --- en promedio se recorre $n/2$ elementos.
    \item \textbf{Espacial}: $O(1)$ --- solo se usan variables auxiliares constantes.
\end{itemize}

\subsection{Casos Límite}
\begin{enumerate}
    \item \textbf{Arreglo vacío} ($n=0$): El bucle no se ejecuta, se retorna $-1$. Correcto.
    \item \textbf{Un solo elemento}: Se compara $A[1]$ con $x$. Si coincide retorna $1$, si no retorna $-1$.
    \item \textbf{Elemento repetido}: Retorna la \emph{primera} aparición (índice menor).
    \item \textbf{Objetivo fuera de rango}: Se recorre todo el arreglo y se retorna $-1$.
    \item \textbf{Todos los elementos iguales al objetivo}: Retorna $1$ (primera posición).
\end{enumerate}

\newpage

% ============================================================
% ALGORITMO 2: CONTAR OCURRENCIAS DE X
% ============================================================
\section{Algoritmo 2: Contar Ocurrencias de \(x\)}

\subsection{Pseudocódigo}
\begin{algorithm}[H]
\caption{Contar Ocurrencias}
\begin{algorithmic}[1]
\Require Arreglo $A[1..n]$, valor $x$
\Ensure Cantidad de veces que $x$ aparece en $A$
\State $c \gets 0$
\For{$i \gets 1$ \textbf{to} $n$}
    \If{$A[i] = x$}
        \State $c \gets c + 1$
    \EndIf
\EndFor
\State \textbf{return} $c$
\end{algorithmic}
\end{algorithm}

\subsection{Invariante de Bucle}
\[
\boxed{\inv(i) \equiv c = |\{k \in \{1, \ldots, i-1\} : A[k] = x\}|}
\]
Es decir, $c$ contiene el número de apariciones de $x$ en las posiciones ya recorridas ($1$ a $i-1$).

\subsection{Demostración}

\subsubsection{Inicialización}
Cuando $i = 1$, el conjunto $\{1, \ldots, 0\} = \emptyset$. Su cardinalidad es $0$ y $c = 0$. Luego $\inv(1)$ es verdadero. $\square$

\subsubsection{Mantenimiento}
Supongamos $\inv(i)$ verdadero. En la iteración $i$:
\begin{itemize}
    \item Si $A[i] = x$: $c$ se incrementa en $1$. Al inicio de la iteración $i+1$, $c = |\{k \leq i : A[k] = x\}|$, que es exactamente $\inv(i+1)$.
    \item Si $A[i] \neq x$: $c$ no cambia. Como $\inv(i)$ contaba correctamente $\{k \leq i-1 : A[k]=x\}$ y $A[i] \neq x$, el conteo para $\{k \leq i : A[k]=x\}$ es el mismo. $\square$
\end{itemize}

\subsubsection{Terminación}
Al terminar el bucle, $i = n+1$. Por el invariante:
\[
c = |\{k \in \{1, \ldots, n\} : A[k] = x\}|
\]
Que es exactamente la cantidad de apariciones de $x$ en $A[1..n]$. $\square$

\subsection{Complejidad}
\begin{itemize}
    \item \textbf{Todos los casos}: $O(n)$ --- siempre se recorre todo el arreglo.
    \item \textbf{Espacial}: $O(1)$.
\end{itemize}

\subsection{Casos Límite}
\begin{enumerate}
    \item \textbf{Arreglo vacío}: $c = 0$. Correcto.
    \item \textbf{$x$ no aparece}: $c = 0$. Correcto.
    \item \textbf{Todos los elementos son $x$}: $c = n$.
    \item \textbf{Un solo elemento}: $c \in \{0, 1\}$ según coincida o no.
\end{enumerate}

\newpage

% ============================================================
% ALGORITMO 3: ENCONTRAR EL MÁXIMO
% ============================================================
\section{Algoritmo 3: Encontrar el Máximo}

\subsection{Pseudocódigo}
\begin{algorithm}[H]
\caption{Encontrar el Máximo}
\begin{algorithmic}[1]
\Require Arreglo $A[1..n]$, $n \geq 1$
\Ensure Elemento máximo de $A$
\State $\text{max} \gets A[1]$
\For{$i \gets 2$ \textbf{to} $n$}
    \If{$A[i] > \text{max}$}
        \State $\text{max} \gets A[i]$
    \EndIf
\EndFor
\State \textbf{return} $\text{max}$
\end{algorithmic}
\end{algorithm}

\subsection{Invariante de Bucle}
\[
\boxed{\inv(i) \equiv \text{max} = \max\{A[1], A[2], \ldots, A[i-1]\}}
\]
El valor de \texttt{max} es el máximo entre los elementos de las posiciones $1$ a $i-1$.

\subsection{Demostración}

\subsubsection{Inicialización}
Cuando $i = 2$: $\text{max} = A[1] = \max\{A[1]\}$. El conjunto $\{A[1]\}$ tiene máximo $A[1]$. $\square$

\subsubsection{Mantenimiento}
Supongamos $\inv(i)$ verdadero: $\text{max} = \max\{A[1], \ldots, A[i-1]\}$.
\begin{itemize}
    \item Si $A[i] > \text{max}$: se actualiza $\text{max} \gets A[i]$. Entonces $\text{max} = \max\{A[1], \ldots, A[i]\}$, que es $\inv(i+1)$.
    \item Si $A[i] \leq \text{max}$: no hay cambio. Como $\text{max} \geq A[i]$ y $\text{max} = \max\{A[1], \ldots, A[i-1]\}$, entonces $\text{max} = \max\{A[1], \ldots, A[i]\}$. $\square$
\end{itemize}

\subsubsection{Terminación}
Cuando el bucle termina, $i = n+1$. Por el invariante:
\[
\text{max} = \max\{A[1], \ldots, A[n]\}
\]
Que es el máximo del arreglo completo. $\square$

\subsection{Complejidad}
\begin{itemize}
    \item \textbf{Todos los casos}: $O(n)$ --- se recorre el arreglo completo siempre.
    \item \textbf{Espacial}: $O(1)$.
\end{itemize}

\subsection{Casos Límite}
\begin{enumerate}
    \item \textbf{Un solo elemento}: El bucle no se ejecuta, retorna $A[1]$.
    \item \textbf{Todos iguales}: $\text{max}$ nunca se actualiza, retorna $A[1]$.
    \item \textbf{Máximo al final}: Se actualiza en la última iteración.
    \item \textbf{Máximo al inicio}: $\text{max}$ nunca se actualiza.
    \item \textbf{Valores negativos}: Funciona correctamente.
\end{enumerate}

\newpage

% ============================================================
% ALGORITMO 4: ENCONTRAR EL MÍNIMO
% ============================================================
\section{Algoritmo 4: Encontrar el Mínimo}

\subsection{Pseudocódigo}
\begin{algorithm}[H]
\caption{Encontrar el Mínimo}
\begin{algorithmic}[1]
\Require Arreglo $A[1..n]$, $n \geq 1$
\Ensure Elemento mínimo de $A$
\State $\text{min} \gets A[1]$
\For{$i \gets 2$ \textbf{to} $n$}
    \If{$A[i] < \text{min}$}
        \State $\text{min} \gets A[i]$
    \EndIf
\EndFor
\State \textbf{return} $\text{min}$
\end{algorithmic}
\end{algorithm}

\subsection{Invariante de Bucle}
\[
\boxed{\inv(i) \equiv \text{min} = \min\{A[1], A[2], \ldots, A[i-1]\}}
\]

\subsection{Demostración}
La demostración es análoga a la del Algoritmo 3, reemplazando $\max$ por $\min$ y la condición $>$ por $<$.

\subsubsection{Inicialización}
$i = 2$: $\text{min} = A[1] = \min\{A[1]\}$. $\square$

\subsubsection{Mantenimiento}
Si $A[i] < \text{min}$, se actualiza y $\text{min} = \min\{A[1],\ldots,A[i]\}$. Si $A[i] \geq \text{min}$, el mínimo no cambia y la propiedad se preserva. $\square$

\subsubsection{Terminación}
$i = n+1$: $\text{min} = \min\{A[1], \ldots, A[n]\}$. $\square$

\subsection{Complejidad}
\begin{itemize}
    \item \textbf{Todos los casos}: $O(n)$.
    \item \textbf{Espacial}: $O(1)$.
\end{itemize}

\subsection{Casos Límite}
\begin{enumerate}
    \item \textbf{Un solo elemento}: Retorna $A[1]$.
    \item \textbf{Todos iguales}: Retorna $A[1]$.
    \item \textbf{Mínimo al inicio}: $\text{min}$ nunca se actualiza.
    \item \textbf{Valores negativos}: Funciona correctamente.
\end{enumerate}

\newpage

% ============================================================
% ALGORITMO 5: VERIFICAR ELEMENTOS POSITIVOS
% ============================================================
\section{Algoritmo 5: Verificar Elementos Positivos}

\subsection{Pseudocódigo}
\begin{algorithm}[H]
\caption{Verificar Todos Positivos}
\begin{algorithmic}[1]
\Require Arreglo $A[1..n]$
\Ensure \textbf{true} si todos los elementos son estrictamente positivos, \textbf{false} en caso contrario
\State $\text{todos\_pos} \gets \textbf{true}$
\For{$i \gets 1$ \textbf{to} $n$}
    \If{$A[i] \leq 0$}
        \State $\text{todos\_pos} \gets \textbf{false}$
        \State \textbf{return} $\textbf{false}$
    \EndIf
\EndFor
\State \textbf{return} $\textbf{true}$
\end{algorithmic}
\end{algorithm}

\subsection{Invariante de Bucle}
\[
\boxed{\inv(i) \equiv \forall k \in \{1, \ldots, i-1\}: A[k] > 0}
\]
Todos los elementos visitados hasta la posición $i-1$ son estrictamente positivos.

\subsection{Demostración}

\subsubsection{Inicialización}
$i = 1$: El conjunto $\{1, \ldots, 0\} = \emptyset$. La afirmación ``$\forall k \in \emptyset: A[k] > 0$'' es \textbf{vacuamente verdadera}. $\square$

\subsubsection{Mantenimiento}
Supongamos $\inv(i)$ verdadero. Si $A[i] \leq 0$, el algoritmo retorna \textbf{false} y no hay siguiente iteración. Si $A[i] > 0$, entonces $\forall k \in \{1, \ldots, i\}: A[k] > 0$, lo que es $\inv(i+1)$. $\square$

\subsubsection{Terminación}
Si el bucle completa todas las iteraciones, $i = n+1$ y $\forall k \in \{1, \ldots, n\}: A[k] > 0$. Se retorna \textbf{true}. Si se retorna \textbf{false} antes, es porque se encontró $A[i] \leq 0$. $\square$

\subsection{Complejidad}
\begin{itemize}
    \item \textbf{Mejor caso}: $O(1)$ --- el primer elemento no es positivo.
    \item \textbf{Peor caso}: $O(n)$ --- todos positivos o el último no lo es.
    \item \textbf{Espacial}: $O(1)$.
\end{itemize}

\subsection{Casos Límite}
\begin{enumerate}
    \item \textbf{Arreglo vacío}: El bucle no se ejecuta, retorna \textbf{true} (vacuamente verdadero).
    \item \textbf{Contiene cero}: Retorna \textbf{false} (se requiere \emph{estrictamente} positivo).
    \item \textbf{Todos negativos}: Retorna \textbf{false} en la primera iteración.
    \item \textbf{Un solo elemento positivo}: Retorna \textbf{true}.
\end{enumerate}

\newpage

% ============================================================
% ALGORITMO 6: VERIFICAR ARREGLO ORDENADO
% ============================================================
\section{Algoritmo 6: Verificar Arreglo Ordenado (No Decreciente)}

\subsection{Pseudocódigo}
\begin{algorithm}[H]
\caption{Verificar Orden No Decreciente}
\begin{algorithmic}[1]
\Require Arreglo $A[1..n]$
\Ensure \textbf{true} si $A[1] \leq A[2] \leq \cdots \leq A[n]$
\State $\text{ordenado} \gets \textbf{true}$
\For{$i \gets 1$ \textbf{to} $n - 1$}
    \If{$A[i] > A[i+1]$}
        \State $\text{ordenado} \gets \textbf{false}$
        \State \textbf{return} $\textbf{false}$
    \EndIf
\EndFor
\State \textbf{return} $\textbf{true}$
\end{algorithmic}
\end{algorithm}

\subsection{Invariante de Bucle}
\[
\boxed{\inv(i) \equiv \forall k \in \{1, \ldots, i-1\}: A[k] \leq A[k+1]}
\]
Todas las parejas consecutivas visitadas hasta la posición $i$ están en orden no decreciente.

\subsection{Demostración}

\subsubsection{Inicialización}
$i = 1$: No se ha verificado ninguna pareja. La afirmación es vacuamente verdadera. $\square$

\subsubsection{Mantenimiento}
Si $A[i] \leq A[i+1]$, la propiedad se extiende a la posición $i$. Si $A[i] > A[i+1]$, se retorna \textbf{false}. En el caso de continuación, $\inv(i+1)$ se mantiene. $\square$

\subsubsection{Terminación}
Si el bucle completa todas las iteraciones ($i = n$), entonces $\forall k \in \{1, \ldots, n-1\}: A[k] \leq A[k+1]$, lo que define un arreglo no decreciente. $\square$

\subsection{Complejidad}
\begin{itemize}
    \item \textbf{Mejor caso}: $O(1)$ --- la primera pareja viola el orden.
    \item \textbf{Peor caso}: $O(n)$ --- el arreglo está ordenado.
    \item \textbf{Espacial}: $O(1)$.
\end{itemize}

\subsection{Casos Límite}
\begin{enumerate}
    \item \textbf{Arreglo vacío o un elemento}: Retorna \textbf{true} (trivialmente ordenado).
    \item \textbf{Todos iguales}: Retorna \textbf{true} (no decreciente incluye empates).
    \item \textbf{Estrictamente decreciente}: Retorna \textbf{false} en la primera iteración.
    \item \textbf{Elementos duplicados}: Son permitidos en orden no decreciente.
\end{enumerate}

\newpage

% ============================================================
% ALGORITMO 7: SUMA DE UN ARREGLO
% ============================================================
\section{Algoritmo 7: Suma de un Arreglo}

\subsection{Pseudocódigo}
\begin{algorithm}[H]
\caption{Suma de Arreglo}
\begin{algorithmic}[1]
\Require Arreglo $A[1..n]$
\Ensure $\sum_{k=1}^{n} A[k]$
\State $s \gets 0$
\For{$i \gets 1$ \textbf{to} $n$}
    \State $s \gets s + A[i]$
\EndFor
\State \textbf{return} $s$
\end{algorithmic}
\end{algorithm}

\subsection{Invariante de Bucle}
\[
\boxed{\inv(i) \equiv s = \sum_{k=1}^{i-1} A[k]}
\]

\subsection{Demostración}

\subsubsection{Inicialización}
$i = 1$: $s = 0 = \sum_{k=1}^{0} A[k] = 0$ (suma vacía). $\square$

\subsubsection{Mantenimiento}
En la iteración $i$: $s \gets s + A[i]$. Si antes $s = \sum_{k=1}^{i-1} A[k]$, después $s = \sum_{k=1}^{i} A[k]$, que es $\inv(i+1)$. $\square$

\subsubsection{Terminación}
$i = n+1$: $s = \sum_{k=1}^{n} A[k]$, que es la suma total. $\square$

\subsection{Complejidad}
\begin{itemize}
    \item \textbf{Todos los casos}: $O(n)$.
    \item \textbf{Espacial}: $O(1)$.
\end{itemize}

\subsection{Casos Límite}
\begin{enumerate}
    \item \textbf{Arreglo vacío}: Retorna $0$.
    \item \textbf{Todos negativos}: La suma es negativa, correcto.
    \item \textbf{Un solo elemento}: Retorna $A[1]$.
    \item \textbf{Overflow}: Discutir límites del tipo de dato.
\end{enumerate}

\newpage

% ============================================================
% ALGORITMO 8: INVERTIR UN ARREGLO
% ============================================================
\section{Algoritmo 8: Invertir un Arreglo}

\subsection{Pseudocódigo}
\begin{algorithm}[H]
\caption{Invertir Arreglo in-place}
\begin{algorithmic}[1]
\Require Arreglo $A[1..n]$
\Ensure $A$ con sus elementos en orden inverso
\For{$i \gets 1$ \textbf{to} $\lfloor n/2 \rfloor$}
    \State $\text{temp} \gets A[i]$
    \State $A[i] \gets A[n - i + 1]$
    \State $A[n - i + 1] \gets \text{temp}$
\EndFor
\end{algorithmic}
\end{algorithm}

\subsection{Invariante de Bucle}
Definimos las posiciones simétricas: para cada $j$, la posición espejo es $n - j + 1$.
\[
\boxed{\inv(i) \equiv \begin{cases}
\forall k \in \{1, \ldots, i-1\}: A[k] = B[n-k+1] & \text{(izquierda: ya trae el valor original de la derecha)}\\
\forall k \in \{n-i+2, \ldots, n\}: A[k] = B[n-k+1] & \text{(derecha: ya trae el valor original de la izquierda)}
\end{cases}}
\]
 donde $B$ es el arreglo original. Alternativamente:
\[
\inv(i) \equiv \text{``Los primeros } i-1 \text{ y los últimos } i-1 \text{ elementos están intercambiados correctamente.''}
\]

\subsection{Demostración}

\subsubsection{Inicialización}
$i = 1$: No se ha intercambiado ningún elemento. El arreglo es igual al original. Verdadero. $\square$

\subsubsection{Mantenimiento}
En la iteración $i$, se intercambian $A[i]$ y $A[n-i+1]$. Si antes los primeros $i-1$ y últimos $i-1$ estaban invertidos, ahora los primeros $i$ y últimos $i$ también lo están. $\square$

\subsubsection{Terminación}
Cuando $i = \lfloor n/2 \rfloor + 1$, todos los pares simétricos $(A[k], A[n-k+1])$ han sido intercambiados para $k = 1, \ldots, \lfloor n/2 \rfloor$. El arreglo está completamente invertido. $\square$

\subsection{Complejidad}
\begin{itemize}
    \item \textbf{Todos los casos}: $O(n)$ --- se realizan $\lfloor n/2 \rfloor$ intercambios.
    \item \textbf{Espacial}: $O(1)$ --- in-place con una variable temporal.
\end{itemize}

\subsection{Casos Límite}
\begin{enumerate}
    \item \textbf{Arreglo vacío}: No se ejecuta ningún intercambio.
    \item \textbf{Un solo elemento}: No se ejecuta ningún intercambio.
    \item \textbf{Par de elementos}: Se intercambian una vez.
    \item \textbf{Número impar de elementos}: El elemento central no se mueve.
    \item \textbf{Palíndromo}: El resultado es igual al original.
\end{enumerate}

\newpage

% ============================================================
% ALGORITMO 9: MÁXIMO SUBARREGLO (KADANE)
% ============================================================
\section{Algoritmo 9: Máximo Subarreglo (Kadane)}

\subsection{Pseudocódigo}
\begin{algorithm}[H]
\caption{Máximo Subarreglo (Kadane)}
\begin{algorithmic}[1]
\Require Arreglo $A[1..n]$
\Ensure Mayor suma de un subarreglo contiguo
\State $\text{max\_actual} \gets A[1]$
\State $\text{max\_global} \gets A[1]$
\For{$i \gets 2$ \textbf{to} $n$}
    \State $\text{max\_actual} \gets \max(A[i],\ \text{max\_actual} + A[i])$
    \If{$\text{max\_actual} > \text{max\_global}$}
        \State $\text{max\_global} \gets \text{max\_actual}$
    \EndIf
\EndFor
\State \textbf{return} $\text{max\_global}$
\end{algorithmic}
\end{algorithm}

\subsection{Invariante de Bucle}
\[
\boxed{\inv(i) \equiv \begin{cases}
\text{max\_global} = \max_{1 \leq j \leq k \leq i} \sum_{t=j}^{k} A[t]\\
\text{max\_actual} = \max_{1 \leq j \leq i} \sum_{t=j}^{i} A[t] \text{ (máximo subarreglo terminando en } i\text{)}
\end{cases}}
\]

\subsection{Demostración}

\subsubsection{Inicialización}
$i = 1$: $\text{max\_actual} = A[1] = \max$ subarreglo terminando en $1$. $\text{max\_global} = A[1] = \max$ subarreglo en $A[1..1]$. $\square$

\subsubsection{Mantenimiento}
En la iteración $i$, el subarreglo máximo terminando en $i$ es:
\begin{itemize}
    \item $A[i]$ solo (reiniciar subarreglo), o
    \item $A[i]$ concatenado al mejor subarreglo terminando en $i-1$.
\end{itemize}
La operación $\text{max\_actual} \gets \max(A[i], \text{max\_actual} + A[i])$ captura exactamente esta decisión. Luego \texttt{max\_global} se actualiza si \texttt{max\_actual} es mayor. $\square$

\subsubsection{Terminación}
Al finalizar, \texttt{max\_global} contiene la máxima suma de cualquier subarreglo contiguo en $A[1..n]$. $\square$

\subsection{Complejidad}
\begin{itemize}
    \item \textbf{Todos los casos}: $O(n)$ --- una pasada por el arreglo.
    \item \textbf{Espacial}: $O(1)$.
\end{itemize}

\subsection{Casos Límite}
\begin{enumerate}
    \item \textbf{Un solo elemento}: Retorna $A[1]$.
    \item \textbf{Todos negativos}: Retorna el menos negativo (el máximo).
    \item \textbf{Todos positivos}: Retorna la suma de todo el arreglo.
    \item \textbf{Mezcla}: El algoritmo maneja correctamente cambios de signo.
\end{enumerate}

\newpage

% ============================================================
% ALGORITMO 10: SUMA MÁXIMA DE VENTANA FIJA
% ============================================================
\section{Algoritmo 10: Suma Máxima de Ventana Fija}

\subsection{Pseudocódigo}
\begin{algorithm}[H]
\caption{Suma Máxima de Ventana Fija (tamaño $k$)}
\begin{algorithmic}[1]
\Require Arreglo $A[1..n]$, tamaño de ventana $k \leq n$
\Ensure Mayor suma entre ventanas contiguas de tamaño $k$
\State $s \gets \sum_{j=1}^{k} A[j]$ \Comment{Suma de la primera ventana}
\State $\text{max\_suma} \gets s$
\For{$i \gets k + 1$ \textbf{to} $n$}
    \State $s \gets s + A[i] - A[i - k]$ \Comment{Deslizar ventana}
    \If{$s > \text{max\_suma}$}
        \State $\text{max\_suma} \gets s$
    \EndIf
\EndFor
\State \textbf{return} $\text{max\_suma}$
\end{algorithmic}
\end{algorithm}

\subsection{Invariante de Bucle}
\[
\boxed{\inv(i) \equiv s = \sum_{j=i-k}^{i-1} A[j] \quad \wedge \quad \text{max\_suma} = \max_{t=k}^{i-1} \sum_{j=t-k+1}^{t} A[j]}
\]
$s$ es la suma de la ventana que termina en la posición $i-1$, y \texttt{max\_suma} es la mayor suma de todas las ventanas ya procesadas.

\subsection{Demostración}

\subsubsection{Inicialización}
Antes del bucle: $s = \sum_{j=1}^{k} A[j]$ (primera ventana). $\text{max\_suma} = s$. Correcto para la primera ventana. $\square$

\subsubsection{Mantenimiento}
En la iteración $i$: la ventana se desliza una posición. Se suma $A[i]$ (nuevo elemento) y se resta $A[i-k]$ (elemento que sale). La nueva suma $s$ corresponde a la ventana $[i-k+1, i]$. Si $s > \text{max\_suma}$, se actualiza. $\square$

\subsubsection{Terminación}
Al finalizar, todas las ventanas de tamaño $k$ fueron evaluadas y \texttt{max\_suma} contiene la mayor suma. $\square$

\subsection{Complejidad}
\begin{itemize}
    \item \textbf{Todos los casos}: $O(n)$ --- cada elemento entra y sale de la ventana una vez.
    \item \textbf{Espacial}: $O(1)$.
\end{itemize}

\subsection{Casos Límite}
\begin{enumerate}
    \item \textbf{$k = n$}: Una sola ventana, retorna la suma de todo el arreglo.
    \item \textbf{$k = 1$}: Retorna el máximo elemento.
    \item \textbf{Todos negativos}: Retorna la suma menos negativa.
    \item \textbf{$n = k$}: Solo hay una ventana.
\end{enumerate}

\newpage

% ============================================================
% ALGORITMO 11: BUBBLE SORT
% ============================================================
\section{Algoritmo 11: Bubble Sort}

\subsection{Pseudocódigo}
\begin{algorithm}[H]
\caption{Bubble Sort}
\begin{algorithmic}[1]
\Require Arreglo $A[1..n]$
\Ensure $A$ ordenado en orden no decreciente
\For{$i \gets 1$ \textbf{to} $n - 1$}
    \For{$j \gets 1$ \textbf{to} $n - i$}
        \If{$A[j] > A[j+1]$}
            \State $\text{intercambiar}(A[j], A[j+1])$
        \EndIf
    \EndFor
\EndFor
\end{algorithmic}
\end{algorithm}

\subsection{Invariante de Bucle}
Para el bucle externo (índice $i$):
\[
\boxed{\inv(i) \equiv \begin{cases}
\text{Los últimos } i-1 \text{ elementos de } A \text{ están en su posición final y son los } i-1 \text{ mayores.}\\
A[j] \leq A[k] \text{ para todo } j > n-i \text{ y } k > j.
\end{cases}}
\]
Para el bucle interno (índice $j$):
\[
\inv_{\text{int}}(j) \equiv A[j] = \max\{A[1], \ldots, A[j]\} \text{ o el mayor ha sido ``burbujeado'' hacia la derecha.}
\]

\subsection{Demostración}

\subsubsection{Inicialización}
$i = 1$: No se han colocado elementos finales. El invariante es trivialmente verdadero. $\square$

\subsubsection{Mantenimiento (Bucle Interno)}
Durante la iteración $j$ del bucle interno, el mayor elemento de $A[1..n-i+1]$ ``burbujea'' hacia la posición $n-i+1$. Después del bucle interno, $A[n-i+1]$ es el elemento en su posición correcta. $\square$

\subsubsection{Mantenimiento (Bucle Externo)}
Después de la iteración $i$, los últimos $i$ elementos están ordenados y en sus posiciones finales. $\square$

\subsubsection{Terminación}
Cuando $i = n$, los últimos $n-1$ elementos están en posición final, lo que implica que todo el arreglo está ordenado. $\square$

\subsection{Complejidad}
\begin{itemize}
    \item \textbf{Mejor caso}: $O(n)$ --- con optimización (bandera \texttt{swapped}), arreglo ya ordenado.
    \item \textbf{Peor caso}: $O(n^2)$ --- arreglo en orden inverso.
    \item \textbf{Promedio}: $O(n^2)$.
    \item \textbf{Espacial}: $O(1)$ --- in-place.
\end{itemize}

\subsection{Casos Límite}
\begin{enumerate}
    \item \textbf{Arreglo vacío o un elemento}: No se ejecuta ningún intercambio.
    \item \textbf{Ya ordenado}: Con optimización, $O(n)$; sin ella, $O(n^2)$.
    \item \textbf{Orden inverso}: Peor caso, cada par se intercambia.
    \item \textbf{Elementos duplicados}: El operador $>$ garantiza estabilidad (no intercambia iguales).
\end{enumerate}

\newpage

% ============================================================
% ALGORITMO 12: INSERTION SORT
% ============================================================
\section{Algoritmo 12: Insertion Sort}

\subsection{Pseudocódigo}
\begin{algorithm}[H]
\caption{Insertion Sort}
\begin{algorithmic}[1]
\Require Arreglo $A[1..n]$
\Ensure $A$ ordenado en orden no decreciente
\For{$i \gets 2$ \textbf{to} $n$}
    \State $\text{clave} \gets A[i]$
    \State $j \gets i - 1$
    \While{$j \geq 1 \wedge A[j] > \text{clave}$}
        \State $A[j+1] \gets A[j]$
        \State $j \gets j - 1$
    \EndWhile
    \State $A[j+1] \gets \text{clave}$
\EndFor
\end{algorithmic}
\end{algorithm}

\subsection{Invariante de Bucle}
Para el bucle externo ($i$):
\[
\boxed{\inv(i) \equiv \text{El subarreglo } A[1..i-1] \text{ está ordenado en orden no decreciente.}}
\]
Para el bucle interno ($j$):
\[
\inv_{\text{int}}(j) \equiv \text{Todos los elementos de } A[j+2..i] \text{ son mayores que } A[j+1] \text{ y el subarreglo se mantiene ordenado.}
\]

\subsection{Demostración}

\subsubsection{Inicialización}
$i = 2$: El subarreglo $A[1..1]$ contiene un solo elemento, que trivialmente está ordenado. $\square$

\subsubsection{Mantenimiento}
Supongamos $A[1..i-1]$ está ordenado. Se inserta $A[i] = \text{clave}$ en su posición correcta:
\begin{itemize}
    \item Los elementos mayores que \texttt{clave} se desplazan a la derecha.
    \item \texttt{clave} se inserta en la posición $j+1$ donde $A[j] \leq \texttt{clave} \leq A[j+1]$.
    \item El subarreglo $A[1..i]$ queda ordenado. $\square$
\end{itemize}

\subsubsection{Terminación}
Cuando el bucle termina, $i = n+1$. El subarreglo $A[1..n]$ está ordenado. $\square$

\subsection{Complejidad}
\begin{itemize}
    \item \textbf{Mejor caso}: $O(n)$ --- arreglo ya ordenado, cada inserción es $O(1)$.
    \item \textbf{Peor caso}: $O(n^2)$ --- arreglo en orden inverso, cada inserción mueve todos los elementos.
    \item \textbf{Promedio}: $O(n^2)$.
    \item \textbf{Espacial}: $O(1)$ --- in-place.
\end{itemize}

\subsection{Casos Límite}
\begin{enumerate}
    \item \textbf{Arreglo vacío o un elemento}: No se ejecuta el bucle.
    \item \textbf{Ya ordenado}: $O(n)$ --- las inserciones son triviales.
    \item \textbf{Orden inverso}: Peor caso, cada elemento se mueve hasta el inicio.
    \item \textbf{Estabilidad}: Insertion Sort es estable (no cambia el orden de elementos iguales).
\end{enumerate}

\newpage

% ============================================================
% ALGORITMO 13: SELECTION SORT
% ============================================================
\section{Algoritmo 13: Selection Sort}

\subsection{Pseudocódigo}
\begin{algorithm}[H]
\caption{Selection Sort}
\begin{algorithmic}[1]
\Require Arreglo $A[1..n]$
\Ensure $A$ ordenado en orden no decreciente
\For{$i \gets 1$ \textbf{to} $n - 1$}
    \State $\text{min\_idx} \gets i$
    \For{$j \gets i + 1$ \textbf{to} $n$}
        \If{$A[j] < A[\text{min\_idx}]$}
            \State $\text{min\_idx} \gets j$
        \EndIf
    \EndFor
    \State $\text{intercambiar}(A[i], A[\text{min\_idx}])$
\EndFor
\end{algorithmic}
\end{algorithm}

\subsection{Invariante de Bucle}
\[
\boxed{\inv(i) \equiv \begin{cases}
A[1..i-1] \text{ está ordenado.}\\
\forall k \in \{1, \ldots, i-1\}, \forall m \in \{i, \ldots, n\}: A[k] \leq A[m]
\end{cases}}
\]
Los primeros $i-1$ elementos están ordenados y son todos menores o iguales a los elementos restantes.

\subsection{Demostración}

\subsubsection{Inicialización}
$i = 1$: El subarreglo $A[1..0]$ es vacío, trivialmente ordenado. La condición de dominancia es vacuamente verdadera. $\square$

\subsubsection{Mantenimiento}
En la iteración $i$:
\begin{enumerate}
    \item Se busca el mínimo en $A[i..n]$ y se obtiene su índice \texttt{min\_idx}.
    \item Se intercambia $A[i]$ con $A[\text{min\_idx}]$.
    \item Ahora $A[i]$ es el menor de $A[i..n]$.
    \item Dado que $A[1..i-1]$ estaba ordenado y todos sus elementos son $\leq A[i]$ (por la búsqueda del mínimo), $A[1..i]$ está ordenado. $\square$
\end{enumerate}

\subsubsection{Terminación}
Cuando $i = n$, $A[1..n-1]$ está ordenado y todos sus elementos son $\leq A[n]$. Por transitividad, $A$ está completamente ordenado. $\square$

\subsection{Complejidad}
\begin{itemize}
    \item \textbf{Todos los casos}: $O(n^2)$ --- siempre se busca el mínimo en el subarreglo restante.
    \item \textbf{Espacial}: $O(1)$ --- in-place.
    \item \textbf{Intercambios}: Como máximo $n-1$ intercambios (ventaja sobre Bubble Sort).
\end{itemize}

\subsection{Casos Límite}
\begin{enumerate}
    \item \textbf{Arreglo vacío o un elemento}: No se ejecuta.
    \item \textbf{Ya ordenado}: Sigue siendo $O(n^2)$, pero solo $n-1$ intercambios.
    \item \textbf{Orden inverso}: $O(n^2)$ y todos los intercambios son necesarios.
    \item \textbf{No estable}: Selection Sort no es estable en general (el intercambio puede cambiar el orden de elementos iguales).
\end{enumerate}

\newpage

% ============================================================
% ALGORITMO 14: MERGE SORT
% ============================================================
\section{Algoritmo 14: Merge Sort}

\subsection{Pseudocódigo}
\begin{algorithm}[H]
\caption{Merge Sort}
\begin{algorithmic}[1]
\Require Arreglo $A[1..n]$
\Ensure $A$ ordenado en orden no decreciente
\Procedure{MergeSort}{$A, l, r$}
    \If{$l < r$}
        \State $m \gets \lfloor (l + r) / 2 \rfloor$
        \State \Call{MergeSort}{$A, l, m$}
        \State \Call{MergeSort}{$A, m+1, r$}
        \State \Call{Merge}{$A, l, m, r$}
    \EndIf
\EndProcedure
\Statex
\Procedure{Merge}{$A, l, m, r$}
    \State Crear $L \gets A[l..m]$, $R \gets A[m+1..r]$
    \State $i \gets 1$, $j \gets 1$, $k \gets l$
    \While{$i \leq |L| \wedge j \leq |R|$}
        \If{$L[i] \leq R[j]$}
            \State $A[k] \gets L[i]$, $i \gets i+1$
        \Else
            \State $A[k] \gets R[j]$, $j \gets j+1$
        \EndIf
        \State $k \gets k + 1$
    \EndWhile
    \State Copiar elementos restantes de $L$ o $R$ en $A$
\EndProcedure
\end{algorithmic}
\end{algorithm}

\subsection{Invariante de Bucle}

Para el procedimiento \texttt{MergeSort} (recursión):
\[
\inv_{\text{merge}} \equiv \text{``Después de } \texttt{MergeSort}(A, l, r)\text{, el subarreglo } A[l..r] \text{ está ordenado.''}
\]

Para el bucle del procedimiento \texttt{Merge}:
\[
\boxed{\inv(i,j) \equiv \begin{cases}
A[l..k-1] \text{ contiene los } k-l \text{ menores elementos de } L[1..i-1] \cup R[1..j-1] \text{ en orden.}\\
$L[i..|L|]$ y $R[j..|R|]$ contienen los elementos no procesados.
\end{cases}}
\]

\subsection{Demostración}

\subsubsection{Inicialización (MergeSort --- Caso base)}
Si $l \geq r$, el subarreglo tiene 0 o 1 elemento, que está trivialmente ordenado. $\square$

\subsubsection{Inicialización (Merge)}
Al inicio del bucle, $k = l$, $i = 1$, $j = 1$. $A[l..k-1] = A[l..l-1] = \emptyset$. Verdadero vacuamente. $\square$

\subsubsection{Mantenimiento (Merge)}
En cada iteración, se selecciona el menor entre $L[i]$ y $R[j]$. Este se coloca en $A[k]$ en su posición correcta relativa. Los elementos restantes conservan su orden. $\square$

\subsubsection{Terminación (Merge)}
Cuando el bucle termina, uno de los arreglos se agotó. Se copian los elementos restantes del otro (que ya están ordenados). $A[l..r]$ contiene todos los elementos en orden. $\square$

\subsubsection{Terminación (MergeSort)}
Por inducción: si \texttt{MergeSort} funciona para arreglos de tamaño $< n$, entonces después de ordenar $A[l..m]$ y $A[m+1..r]$, \texttt{Merge} combina dos arreglos ordenados en uno ordenado. $\square$

\subsection{Complejidad}
\begin{itemize}
    \item \textbf{Mejor caso}: $O(n \log n)$.
    \item \textbf{Peor caso}: $O(n \log n)$ --- siempre divide en dos mitades.
    \item \textbf{Promedio}: $O(n \log n)$.
    \item \textbf{Espacial}: $O(n)$ --- se requieren arreglos auxiliares para \texttt{Merge}.
\end{itemize}

\subsection{Casos Límite}
\begin{enumerate}
    \item \textbf{Arreglo vacío}: La recursión no se ejecuta.
    \item \textbf{Un solo elemento}: Caso base, no se hace nada.
    \item \textbf{Ya ordenado}: Sigue siendo $O(n \log n)$.
    \item \textbf{Todos iguales}: Sigue siendo $O(n \log n)$.
    \item \textbf{Estabilidad}: Merge Sort es estable si la fusión preserva el orden de elementos iguales ($L[i] \leq R[j]$ en lugar de $<$).
\end{enumerate}

\newpage

% ============================================================
% ALGORITMO 15: QUICK SORT
% ============================================================
\section{Algoritmo 15: Quick Sort}

\subsection{Pseudocódigo}
\begin{algorithm}[H]
\caption{Quick Sort}
\begin{algorithmic}[1]
\Require Arreglo $A[1..n]$
\Ensure $A$ ordenado en orden no decreciente
\Procedure{QuickSort}{$A, l, r$}
    \If{$l < r$}
        \State $p \gets \texttt{Partition}(A, l, r)$
        \State \Call{QuickSort}{$A, l, p-1$}
        \State \Call{QuickSort}{$A, p+1, r$}
    \EndIf
\EndProcedure
\Statex
\Procedure{Partition}{$A, l, r$}
    \State $\text{pivote} \gets A[r]$
    \State $i \gets l - 1$
    \For{$j \gets l$ \textbf{to} $r - 1$}
        \If{$A[j] \leq \text{pivote}$}
            \State $i \gets i + 1$
            \State $\text{intercambiar}(A[i], A[j])$
        \EndIf
    \EndFor
    \State $\text{intercambiar}(A[i+1], A[r])$
    \State \textbf{return} $i + 1$
\EndProcedure
\end{algorithmic}
\end{algorithm}

\subsection{Invariante de Bucle}

Para el procedimiento \texttt{Partition} (bucle con índice $j$):
\[
\boxed{\inv(j) \equiv \begin{cases}
\forall k \in \{l, \ldots, i\}: A[k] \leq \text{pivote}\\
\forall k \in \{i+1, \ldots, j-1\}: A[k] > \text{pivote}
\end{cases}}
\]
Todos los elementos a la izquierda de $i$ son $\leq$ pivote, y todos entre $i+1$ y $j-1$ son $>$ pivote.

\subsection{Demostración}

\subsubsection{Inicialización}
$j = l$: $i = l - 1$. Los rangos $\{l, \ldots, l-1\}$ y $\{l, \ldots, l-1\}$ son vacíos. Verdadero vacuamente. $\square$

\subsubsection{Mantenimiento}
En la iteración $j$:
\begin{itemize}
    \item Si $A[j] \leq \text{pivote}$: Se incrementa $i$ e intercambia $A[i]$ con $A[j]$. Ahora $A[i] \leq \text{pivote}$ y la invariant se mantiene.
    \item Si $A[j] > \text{pivote}$: No se hace nada. $A[j]$ queda en la región de elementos $>$ pivote. $\square$
\end{itemize}

\subsubsection{Terminación}
Al finalizar el bucle, se intercambia $A[i+1]$ con $A[r]$ (el pivote). Ahora:
\begin{itemize}
    \item $A[l..i] \leq \text{pivote}$
    \item $A[i+1] = \text{pivote}$
    \item $A[i+2..r] > \text{pivote}$
\end{itemize}
El pivote está en su posición final $p = i+1$. Las llamadas recursivas ordenan los subarreglos izquierdo y derecho. $\square$

\subsection{Complejidad}
\begin{itemize}
    \item \textbf{Mejor caso}: $O(n \log n)$ --- partición equilibrada.
    \item \textbf{Peor caso}: $O(n^2)$ --- arreglo ya ordenado o inverso ordenado con pivote en extremo.
    \item \textbf{Promedio}: $O(n \log n)$.
    \item \textbf{Espacial}: $O(\log n)$ --- profundidad promedio de la pila de recursión.
\end{itemize}

\subsection{Casos Límite}
\begin{enumerate}
    \item \textbf{Arreglo vacío o un elemento}: No se ejecuta partición.
    \item \textbf{Todos iguales}: Sin optimización, peor caso $O(n^2)$. Con \texttt{Dutch National Flag} o pivote aleatorio, $O(n)$.
    \item \textbf{Ya ordenado}: Peor caso con pivote en el extremo. Mejorar con pivote aleatorio.
    \item \textbf{No estable}: Quick Sort no es estable en su implementación in-place.
    \item \textbf{Pivote aleatorio}: Seleccionar un pivote aleatorio evita el peor caso determinístico.
\end{enumerate}

\newpage

% ============================================================
% ALGORITMO 16: TWO-SUM
% ============================================================
\section{Algoritmo 16: Two-Sum (Par con Suma Objetivo)}

\subsection{Pseudocódigo}
\begin{algorithm}[H]
\caption{Two-Sum (con tabla hash)}
\begin{algorithmic}[1]
\Require Arreglo $A[1..n]$, objetivo $x$
\Ensure Índices $(p, q)$ con $p < q$ tales que $A[p] + A[q] = x$, o $-1$ si no existe
\State $H \gets \emptyset$ \Comment{Tabla hash: valor $\to$ índice}
\For{$i \gets 1$ \textbf{to} $n$}
    \If{$x - A[i] \in H$}
        \State \textbf{return} $(H[x - A[i]],\ i)$
    \EndIf
    \State $H[A[i]] \gets i$
\EndFor
\State \textbf{return} $-1$
\end{algorithmic}
\end{algorithm}

\subsection{Invariante de Bucle}
\[
\boxed{\inv(i) \equiv H = \{(A[k], k) : 1 \leq k \leq i-1\} \ \wedge \ \forall\, j,k \in \{1, \ldots, i-1\}, j < k: A[j] + A[k] \neq x}
\]
La tabla $H$ contiene exactamente los elementos ya procesados (con su índice), y ningún par de entre esos elementos suma el objetivo $x$.

\subsection{Demostración}

\subsubsection{Inicialización}
Cuando $i = 1$, $H = \emptyset$ y el conjunto de pares procesados es vacío. Ambas condiciones son \textbf{vacuamente verdaderas}. $\square$

\subsubsection{Mantenimiento}
Supongamos $\inv(i)$ verdadero al inicio de la iteración $i$. Se inspecciona $A[i]$:
\begin{itemize}
    \item Si $x - A[i] \in H$: existe $k = H[x - A[i]] < i$ con $A[k] = x - A[i]$, luego $A[k] + A[i] = x$. El algoritmo retorna $(k, i)$, que es correcto. No hay iteración siguiente, por lo que el invariante se mantiene trivialmente.
    \item Si $x - A[i] \notin H$: para todo $k < i$ se tiene $A[k] + A[i] \neq x$. Como $\inv(i)$ afirma que ningún par de $\{1, \ldots, i-1\}$ suma $x$, ningún par de $\{1, \ldots, i\}$ suma $x$. Tras insertar $H[A[i]] \gets i$, la tabla contiene los elementos de $\{1, \ldots, i\}$, que es $\inv(i+1)$. $\square$
\end{itemize}

\subsubsection{Terminación}
El bucle termina en dos casos:
\begin{enumerate}
    \item \textbf{Retorno anticipado}: se encontró un par $(k, i)$ con $A[k] + A[i] = x$, lo cual es correcto.
    \item \textbf{Bucle completa todas las iteraciones}: $\inv(n+1)$ afirma que ningún par de $A[1..n]$ suma $x$. Se retorna $-1$, lo cual es correcto. $\square$
\end{enumerate}
Nótese que si existe un par solución $(p, q)$ con $p < q$, en la iteración $q$ la tabla ya contiene $A[p]$ (pues $p < q$), por lo que el algoritmo necesariamente lo encuentra.

\subsection{Complejidad}
\begin{itemize}
    \item \textbf{Todos los casos}: $O(n)$ esperado --- una pasada por el arreglo con operaciones de tabla hash $O(1)$ esperado.
    \item \textbf{Espacial}: $O(n)$ --- la tabla hash almacena hasta $n$ entradas.
\end{itemize}

\subsection{Casos Límite}
\begin{enumerate}
    \item \textbf{Arreglo vacío o de un elemento}: El bucle no encuentra pares y retorna $-1$.
    \item \textbf{Duplicados}: Si $A = [3, 3]$ con $x = 6$, en $i = 2$ se encuentra $3 \in H$ y se retorna $(1, 2)$. El algoritmo nunca usa el mismo elemento dos veces, pues la verificación precede a la inserción.
    \item \textbf{Objetivo $x = 0$}: Un par de opuestos $A = [-3, 2, 3]$ con $x = 0$ retorna $(1, 3)$.
    \item \textbf{Valores negativos}: El complemento $x - A[i]$ puede ser negativo y se maneja correctamente con la tabla hash.
    \item \textbf{Múltiples pares solución}: Retorna el primero que se encuentra al recorrer el arreglo.
\end{enumerate}

\newpage

% ============================================================
% RESUMEN DE COMPLEJIDADES
% ============================================================
\section{Resumen de Complejidades}

\begin{table}[H]
\centering
\begin{tabular}{@{}llll@{}}
\toprule
\textbf{Algoritmo} & \textbf{Mejor} & \textbf{Peor} & \textbf{Promedio} \\
\midrule
1. Búsqueda Lineal & $O(1)$ & $O(n)$ & $O(n)$ \\
2. Contar Ocurrencias & $O(n)$ & $O(n)$ & $O(n)$ \\
3. Encontrar el Máximo & $O(n)$ & $O(n)$ & $O(n)$ \\
4. Encontrar el Mínimo & $O(n)$ & $O(n)$ & $O(n)$ \\
5. Verificar Positivos & $O(1)$ & $O(n)$ & $O(n)$ \\
6. Verificar Ordenado & $O(1)$ & $O(n)$ & $O(n)$ \\
7. Suma de Arreglo & $O(n)$ & $O(n)$ & $O(n)$ \\
8. Invertir Arreglo & $O(n)$ & $O(n)$ & $O(n)$ \\
9. Máx. Subarreglo (Kadane) & $O(n)$ & $O(n)$ & $O(n)$ \\
10. Suma Máx. Ventana Fija & $O(n)$ & $O(n)$ & $O(n)$ \\
11. Bubble Sort & $O(n)^*$ & $O(n^2)$ & $O(n^2)$ \\
12. Insertion Sort & $O(n)$ & $O(n^2)$ & $O(n^2)$ \\
13. Selection Sort & $O(n^2)$ & $O(n^2)$ & $O(n^2)$ \\
14. Merge Sort & $O(n \log n)$ & $O(n \log n)$ & $O(n \log n)$ \\
15. Quick Sort & $O(n \log n)$ & $O(n^2)$ & $O(n \log n)$ \\
16. Two-Sum & $O(n)$ & $O(n)$ & $O(n)$ \\
\bottomrule
\end{tabular}
\caption{Complejidades temporales. (*) Con optimización de bandera.}
\end{table}

\begin{table}[H]
\centering
\begin{tabular}{@{}lll@{}}
\toprule
\textbf{Algoritmo} & \textbf{Espacial} & \textbf{Estable} \\
\midrule
1--10 & $O(1)$ & N/A \\
11. Bubble Sort & $O(1)$ & Sí \\
12. Insertion Sort & $O(1)$ & Sí \\
13. Selection Sort & $O(1)$ & No \\
14. Merge Sort & $O(n)$ & Sí \\
15. Quick Sort & $O(\log n)$ & No \\
16. Two-Sum & $O(n)$ & N/A \\
\bottomrule
\end{tabular}
\caption{Complejidad espacial y estabilidad.}
\end{table}

\end{document}
